\chapter{Erklärung zum Einsatz von KI-gestützten Werkzeugen}

Im Rahmen der Erstellung dieser Bachelorarbeit wurden KI-gestützte Werkzeuge
eingesetzt. Die nachfolgende Tabelle gibt einen Überblick über die verwendeten
Systeme sowie deren jeweiligen Einsatzbereich.

\begin{table}[htbp]
    \centering
  \begin{tabular}{p{3.5cm} p{3.5cm} p{5cm}}
    \toprule
    \textbf{Werkzeug} & \textbf{Anbieter} & \textbf{Verwendungszweck} \\
    \midrule
    Claude Sonnet 4.6 & Anthropic & Unterstützung bei der sprachlichen
                                   Formulierung von Textpassagen sowie bei der Entwicklung und Überprüfung von Quellcode \\
    \addlinespace
    Perplexity       & Perplexity AI & Recherche und Identifikation
                                       wissenschaftlicher Quellen \\
    \addlinespace
    DeepL             & DeepL SE & Übersetzung von Textpassagen und
                                   Fachbegriffen \\
                                   \addlinespace
    v0               & Vercel & KI-gestütztes Prototyping zur Erstellung 
                             der Mockups der Benutzeroberfläche \\
    \bottomrule
  \end{tabular}
  \caption{Eingesetzte KI-gestützte Werkzeuge}
  \label{tab:ki-werkzeuge}
\end{table}

Sämtliche durch KI-Werkzeuge generierten oder unterstützten Inhalte wurden
kritisch geprüft, überarbeitet und eigenverantwortlich in die Arbeit integriert.
Die inhaltliche Verantwortung für alle Aussagen dieser Arbeit liegt
ausschließlich bei dem Verfasser.